\documentclass{beamer}

\usepackage[utf8]{inputenc}
\usepackage{amssymb, amsmath}
\usepackage{amsthm}
\usepackage{comment}

\usepackage{graphicx}
\graphicspath{{./images/}}

\usetheme{metropolis}
\setbeamercolor{block title}{bg=red!20}

\usepackage{smartdiagram}
%%%%%%%%%%%%%%%%%%%%%%%%%%%%%%%%%%%%%%%%%%%%%%%%%%%%%%%%%%%%%%%
% Declaration for argmax
\DeclareMathOperator*{\argmax}{arg\,max}
\DeclareMathOperator*{\argmin}{arg\,min}

\title{English Word Suggestion Based on Part of Speech Ngram}
\subtitle{Hamana Laboratory, Gunma University}
\author{Borann Chanrathnak}

%%%%%%%%%%%%%%%%%%%%%%%%%%%%%%%%%%%%%%%%%%%%%%%%%%%%%%%%%%%%%%%
\begin{document}
\maketitle

%%%%%%%%%%%%%%%%%%%%%%%%%%%%%%%%%%%%%%%%%%%%%%%%%%%%%%%%%%%%%%%

\begin{frame}
\frametitle{Table of Contents}
\tableofcontents
\end{frame}

%%%%%%%%%%%%%%%%%%%%%%%%%%%%%%%%%%%%%%%%%%%%%%%%%%%%%%%%%%%%%%%
\section{Motivation}
\begin{frame}{Motivation}
    \begin{itemize}
        \item Word suggestion for note-taking app on mobile devices
    \end{itemize}
    \centering
    \only<1>{\includegraphics[width=5.5cm]{one}}
    \only<2>{\includegraphics[width=5.5cm]{two}}
    \only<3>{\includegraphics[width=5.5cm]{three}}
\end{frame}

%%%%%%%%%%%%%%%%%%%%%%%%%%%%%%%%%%%%%%%%%%%%%%%%%%%%%%%%%%%%%%%
\begin{comment}
\begin{frame}{How to \textcolor[rgb]{0.8,0.2,0}{Predict}?}
    \centering

    "I am japanese, so I speak" $\rightarrow$ \visible<2->{"japanese"}
\end{frame}
\end{comment}
%%%%%%%%%%%%%%%%%%%%%%%%%%%%%%%%%%%%%%%%%%%%%%%%%%%%%%%%%%%%%%%
\section{Part Of Speech Ngram (POS-Ngram)}
\subsection{Definitions}

\begin{frame}{What does POS-Ngram mean?}
    \visible<1->{
        \textbf{N-gram} 
        is a contiguous sequence of n items from a given sample of text or speech\\
            \textbf{Ex: } I study english
            \begin{itemize}
                \item \textbf{unigram (1-gram)} : (I,), (study,), (english,)
                \item \textbf{bigram (2-gram)} : (I, study), (study, english) 
                \item \textbf{trigram (3-gram)} : (I, study, english)
            \end{itemize}
    }

    \visible<2->{
        \textbf{N-gram model}
        is one of statistical langauge models for predicting the next item (word) based on Markov assumption, and usually abbreviated as N-gram.
    }

    \visible<3->{
        \textbf{POS-Ngram}
            is an improved model using part of speech as a class indicator.
    }
\end{frame}

%%%%%%%%%%%%%%%%%%%%%%%%%%%%%%%%%%%%%%%%%%%%%%%%%%%%%%%%%%%%%%%
\subsection{N-gram}

\begin{frame}{How To Compute Probability of a Sentence?}
How can we compute the joint probability of a sentence?\\
\textit{Ex:} "an apple is on the table" \\
    $$P(an\ apple\ is\ on\ the\ table) = \text{?}$$

\end{frame}
%%%%%%%%%%%%%%%%%%%%%%%%%%%%%%%%%%%%%%%%%%%%%%%%%%%%%%%%%%%%%%%
\begin{frame}{Notation}
    \begin{block}{\textcolor[rgb]{0.8,0.2,0}{Notation}}
        \visible<1->{\begin{itemize}
            \visible<2->{\item To represent the probability of a particular random variable $X_i$ taking on the value "the", or $P(X_i="the")$ we will use simplification $P(the)$.}
            \visible<3->{\item We represent a sequence of N words either as $w_1\cdots w_n$ or $w_1^n$}
            \visible<4->{\item \textbf{Ex:} "an apple is on the table" \\$\rightarrow$ $w_1="an"$, $w_2="apple"$, $w_3="is", \cdots$\\
                \visible<5->{$\rightarrow w_1^4="an\ apple\ is\ on"$}}
        \end{itemize}}
    \end{block}

\end{frame}
%%%%%%%%%%%%%%%%%%%%%%%%%%%%%%%%%%%%%%%%%%%%%%%%%%%%%%%%%%%%%%%

\begin{frame}{Chain Rule Of Probability}
    \visible<1->{\begin{block}{In General}
        $$P(X_1\cdots X_n) = P(X_1)P(X_2\mid X_1)P(X_3\mid X_1^2)\cdots P(X_n\mid X_1^{n-1})$$
    \end{block}}

    \visible<2->{\begin{block}{By applying chain rule to a sequence of words}
        $$P(w_1\cdots w_n) = P(w_1)P(w_2\mid w_1)P(w_3\mid w_1^2)\cdots P(w_n\mid w_1^{n-1})$$
    \end{block}}

    \visible<3->{\begin{block}{Examples}
        \begin{itemize}
            \item $P(an\ apple) = P(an)\times P(apple\mid an)$
            \item $P(an\ apple\ is) = P(an)\times P(apple\mid an) \times P(is\mid an\ apple)$
        \end{itemize}
    \end{block}}

\end{frame}
%%%%%%%%%%%%%%%%%%%%%%%%%%%%%%%%%%%%%%%%%%%%%%%%%%%%%%%%%%%%%%%

\begin{frame}{Let's Apply Chain Rule}
By applying chain rule to the phrase $"an\ apple\ is\ on\ the"$:
    \begin{align*}
        P("an\ apple\ is\ on\ the") &= P(an)\times P(apple\mid an) \times P(is\mid an\ apple)\\
                                  &\times P(on\mid an\ apple\ is) \times P(the\mid an\ apple\ is\ on) 
    \end{align*}
\end{frame}
%%%%%%%%%%%%%%%%%%%%%%%%%%%%%%%%%%%%%%%%%%%%%%%%%%%%%%%%%%%%%%%

\begin{frame}{In practice chain rule does not help}
    \begin{itemize}[<+->]
        \item Language is creative, and any particular context might have never occured before.
        \item We don't know the way to compute the exact probability of a word given a long sequence of preceding words, $P(w_n\mid w_1^{n-1})$.
    \end{itemize}

    %%%%%%%%%%%%%%%%%%%%%%%%%%%%%%%%%%%%%%%%%%%%%%%%%%%%%%%%%%
    \visible<3->{
    \textbf{Ex:}
    \begin{align*}
        P("an\ apple\ is\ on\ the") &= P(an)\times P(apple\mid an) \times P(is\mid an\ apple)\\
                                  &\times P(on\mid an\ apple\ is) \times 
                                  \only<3>{P(the\mid an\ apple\ is\ on)}
                                  \only<4>{\underline{P(the\mid an\ apple\ is\ on)}}
    \end{align*}}
    %%%%%%%%%%%%%%%%%%%%%%%%%%%%%%%%%%%%%%%%%%%%%%%%%%%%%%%%%%

\end{frame}
%%%%%%%%%%%%%%%%%%%%%%%%%%%%%%%%%%%%%%%%%%%%%%%%%%%%%%%%%%%%%%%

\begin{frame}{The Presence of N-gram}
    The general equation for this n-gram approximation to the conditional probability of the next word in a sequence is
    $$P(w_n\mid w_1^{n-1}) \approx P(w_n\mid w_{n-N+1}^{n-1})$$

    \visible<2->{\begin{block}{In case of Bigram (2-gram)}
        When we use bigram model to predict the conditional probability of the next word, we can approximate by
        $$P(w_n\mid w_1^{n-1}) \approx P(w_n\mid w_{n-1})$$
    \end{block}}

    \visible<3->{\begin{block}{Example}
        $$P(the\mid \ an\ apple\ is\ on) \approx P(the\mid \ on)$$
    \end{block}}
\end{frame}

%%%%%%%%%%%%%%%%%%%%%%%%%%%%%%%%%%%%%%%%%%%%%%%%%%%%%%%%%%%%%%%

\begin{frame}{Bigram in Practice}
By supposing that C is the counts of word from a corpus.
    $$P(w_n\mid w_{n-1}) = \frac{C(w_{n-1}w_n)}{\sum_{w\in dict}C(w_{n-1}w)}$$

    \visible<2->{\begin{block}{Example}
        $$P(the\mid \ on) = \frac{C(on\ the)}{\sum_{x}C(on\ x)}$$
    \end{block}}
    \visible<3->{(\textit{on x}) means (on the), (on board), (on time), (on that), $\cdots$}

\end{frame}

%%%%%%%%%%%%%%%%%%%%%%%%%%%%%%%%%%%%%%%%%%%%%%%%%%%%%%%%%%%%%%%

\begin{frame}{The probability of a complete word sequence}
    Given the \textbf{Bigram} assumption for the probability of an individual word, we can compute the probability of complete word sequence by:
    $$P(w_1^n) \approx \prod_{k=1}^nP(w_k\mid w_{k-1})$$

    % Compare Chain Rule with Bigram based Calculation
    \begin{block}{\only<1>{Chain Rule}\onslide<2>{Bigram Assumption}}
        % CHAIN RULE BASED CALCULATION
        \only<1>{
            \begin{align*}
            P("an\ apple\ is\ on\ the") &= P(an)\times P(apple\mid an) \times P(is\mid an\ apple)\\
                                      &\times P(on\mid an\ apple\ is) \times P(the\mid an\ apple\ is\ on) 
        \end{align*}}
        % BIGRAM BASED CALCULATION
        \onslide<2>{
        \begin{align*}
        P("an\ apple\ is\ on\ the") &= P(apple\mid an) \times P(is\mid apple) \\
                                    &\times P(on\mid is) \times P(the\mid on)
        \end{align*}}
    \end{block}

\end{frame}

%%%%%%%%%%%%%%%%%%%%%%%%%%%%%%%%%%%%%%%%%%%%%%%%%%%%%%%%%%%%%%%

\begin{frame}{OOV \& Smoothing}
    \textbf{OOV} stands for \textbf{Out Of Vocabulary}\\
    \begin{itemize}
        \item Words appear only in a \textcolor{purple}{test set} but not in the \textcolor{purple}{training set}.
        \item OOV problem occurs even when we work on big data.
    \end{itemize}
\end{frame}
%%%%%%%%%%%%%%%%%%%%%%%%%%%%%%%%%%%%%%%%%%%%%%%%%%%%%%%%%%%%%%%

\begin{frame}{Types of smoothing}
    \begin{itemize}
        \item Laplace smoothing (Add-one smoothing)
        \item Add-k smoothing
        \item \only<1>{Interpolation}
              \onslide<2->{\textcolor{purple}{Interpolation}}
        \item $\cdots$
    \end{itemize}
    \onslide<3->{
    \begin{block}{Interpolation}
        We mix the probability estimates from all the n-gram estimators, weighing and combining the trigram, bigram, and unigram counts.}
    \end{block}
\end{frame}
%%%%%%%%%%%%%%%%%%%%%%%%%%%%%%%%%%%%%%%%%%%%%%%%%%%%%%%%%%%%%%%

\begin{frame}{Interpolation}

    \begin{block}{Bigram + Unigram}
        $$\hat{P}(table\mid the) = \lambda_1P(table\mid the) + \lambda_2P(table)$$
    \end{block}

    \visible<2->{\begin{block}{Trigram + Bigram + Unigram}
        \begin{align*}
            \hat{P}(table\mid on\ the) &= \lambda_1P(table\mid on\ the)\\
                                         &+ \lambda_2P(table\mid the)\\
                                         &+ \lambda_3P(table)
        \end{align*}
    \end{block}}

    \visible<3->{where we choose $\lambda_i$ such that $\sum_i\lambda_i = 1$}
\end{frame}
%%%%%%%%%%%%%%%%%%%%%%%%%%%%%%%%%%%%%%%%%%%%%%%%%%%%%%%%%%%%%%%

\begin{frame}{How are the $\lambda$ values set?}
    \begin{itemize}
        \item They can be learned from a \textcolor{purple}{held-out} corpus
        \item Can be found by \textcolor{purple}{EM} algorithm
        \item For the purpose of this project, we assume without loss of generality that $\lambda_i > \lambda_j (\forall i < j)$
    \end{itemize}
\end{frame}
%%%%%%%%%%%%%%%%%%%%%%%%%%%%%%%%%%%%%%%%%%%%%%%%%%%%%%%%%%%%%%%

\subsection{POS-Ngram}
\begin{frame}{POS-Ngram}

    \begin{block}{Formula}
        $$P(w_n\mid w_1^{n-1}) = \sum_{c_n}P(w_n\mid c_n)\times P(c_n\mid c_{n-N+1}^{n-1})$$
        where $c_n$ is a class of $w_n$\\
        (In this context, $c_n$ is considered a part of speech of the word $w_n$)
    \end{block}

    \visible<2->{\begin{block}{Example}
        \centering
        \begin{tabular}{|l|l|}
            \hline
            $c_1$:Noun & cat, dog, thought, $\cdots$\\\hline
            $c_2$:Verb & go, speak, $\cdots$\\\hline
            $c_1$:Noun, $c_2$:Verb & play, $\cdots$\\\hline
        \end{tabular}\\
    \end{block}}

    \flushleft
    
    \visible<3->{It means that to predict the next word $w_n$
    \begin{enumerate}
        \item $P(w_n\mid c_n)=\frac{C(w_n,c_n)}{C(c_n)}$
        \item Compute $P(c_n\mid c_{n-N+1}^{n-1}) \onslide<4->{\sim P(w_n\mid w_{n-N+1}^{n-1})$}
        \item $w_n^\ast$ is determined by $\argmax_{w_n} P(w_n\mid w_1^{n-1})$
    \end{enumerate}}
\end{frame}

%%%%%%%%%%%%%%%%%%%%%%%%%%%%%%%%%%%%%%%%%%%%%%%%%%%%%%%%%%%%%%%
\begin{comment}
\begin{frame}{POS-Ngram}
    \begin{block}{Example}
        \begin{itemize}
            \item $"I\ study\ english"$
            \item $Pronoun\ Verb\ Noun"$
            \item $Pronoun\ Noun\ Adjective"$
            \item $\cdots$
        \end{itemize}
    \end{block}
\end{frame}
\end{comment}
%%%%%%%%%%%%%%%%%%%%%%%%%%%%%%%%%%%%%%%%%%%%%%%%%%%%%%%%%%%%%%%
\section{Implementation}
%%%%%%%%%%%%%%%%%%%%%%%%%%%%%%%%%%%%%%%%%%%%%%%%%%%%%%%%%%%%%%%
\subsection{Programming Structure}
\begin{frame}{Programming Structure}
    \smartdiagramset{}
    \begin{center}
        \smartdiagram[priority descriptive diagram]{PosNgram.py,Predict.py,Ui.py}
    \end{center}
\end{frame}
%%%%%%%%%%%%%%%%%%%%%%%%%%%%%%%%%%%%%%%%%%%%%%%%%%%%%%%%%%%%%%%

\subsection{Tools}
\begin{frame}{Implementation}
    \begin{itemize}
        \item Programming Language : Python
        \item Ui : Tkinter
        \item NLTK (Natural Language ToolKit)
    \end{itemize}
\end{frame}

%%%%%%%%%%%%%%%%%%%%%%%%%%%%%%%%%%%%%%%%%%%%%%%%%%%%%%%%%%%%%%%

\begin{frame}{References}
    \begin{enumerate}
        \item Speech and Language Processing (Chapter 4) (Daniel Jurafsky - Stanford University, 1999)
        \item N-gram Language Modeling Tutorial (Lecture notes courtesy of Prof. Mari Ostendorf, 2007-06-21)
        \item Probabilistic Language Model (Chapter 3) (Kenji KITA, University of Tokyo Press, 1999)
    \end{enumerate}
\end{frame}
%%%%%%%%%%%%%%%%%%%%%%%%%%%%%%%%%%%%%%%%%%%%%%%%%%%%%%%%%%%%%%%
\begin{frame}
    \centering
    \Huge DEMO
\end{frame}

%%%%%%%%%%%%%%%%%%%%%%%%%%%%%%%%%%%%%%%%%%%%%%%%%%%%%%%%%%%%%%%
\end{document}
